%!TEX program = xelatex
\documentclass[12pt,a4paper]{ctexart}

\usepackage[left=1.8cm,right=1.8cm,top=1.4cm,bottom=1.2cm]{geometry}
\usepackage{setspace}
\setstretch{1.0}
\usepackage{booktabs}
\usepackage[colorlinks=true,linkcolor=blue,urlcolor=blue]{hyperref}
\usepackage{titlesec}
\titleformat{\section}{\large\bfseries}{}{0em}{}
\titlespacing{\section}{0pt}{2pt}{0pt}
\titleformat{\subsection}{\normalsize\bfseries}{}{0em}{}
\titlespacing{\subsection}{0pt}{1pt}{0pt}

\begin{document}

\title{\vspace{-1.2cm}\textbf{面向生物医学文献的多策略多标签\\[1pt]分类方法系统比较研究——以空间转录组学文献为例}}
\author{李哲夫\quad 2023011400\quad 生32\\[1pt]
\texttt{\footnotesize\url{https://github.com/zf-li23/pubmed-spatial-tracker}}}
\date{}
\maketitle
\vspace{-1.0cm}

\section{一、选题背景与意义}

生物医学文献正以指数速度增长——空间转录组学领域PubMed年发文量从2020年之前不足500篇激增至2025年4,150篇。自动文献分类面临三重挑战：(1)~文本稀疏，词袋模型无法捕捉医学术语语义关联；(2)~标签空间复杂，一篇文献同时涉及研究类别、生物领域、技术平台和分析方法，构成多标签问题；(3)~标注成本极高，大规模人工标注不现实。

本课题在\textbf{四个独立数据集}上系统比较多策略机器学习方法：(1)~OHSUMED\footnote{\url{https://trec.nist.gov/data/filtering/t9.filtering.tar.gz}}（$\sim$294K篇，14,466个MeSH词），(2)~PubMed-MultiLabel\footnote{\url{https://www.kaggle.com/datasets/owaiskhan9654/pubmed-multilabel-text-classification}}（10K篇，15个MeSH顶级类别），(3)~PGB\footnote{Lee \& Ho, arXiv:2305.02691, 2023}（$\sim$30M篇，含MeSH层级+引用图），(4)~自建空间转录组学文献库（约数千篇，待通过PubMed检索构建）。四个数据集从标签粒度（15$\to$14,466类）、信息维度（纯文本$\to$异构图）和标注完备性（全标注$\to$待标注）形成天然梯度，支撑贝叶斯学习、集成学习、深度学习、无监督学习、基于实例的学习、图表示学习等模块的算法系统比较。

本次为\textbf{修订重提交}。相比初版，我深入研究了生物医学文献分类领域，理解了MeSH主题词层级体系，重构了项目架构，新增两个高质量全标注基准数据集，并找到利用DeepSeek API（OpenAI兼容格式，\texttt{deepseek-v4}模型）进行批量标注\footnote{Touchent et al., arXiv:2506.20331, 2025}的方法来解决目标数据集标注问题。算法方面，受PGB启发新增了\textbf{图表示学习方法}（node2vec、GCN、GraphSAGE）。为严格遵守课程AI使用规定，我已删除仓库全部AI生成代码，将从头手写并全程使用Git记录。本项目不属于任何课程或SRT项目。

\section{二、拟采用的特征与算法}

\subsection*{2.1 文本表示}

采用4种互补的文本表示方法，从不同粒度捕捉文献语义：
\begin{itemize}
  \setlength{\itemsep}{-2pt}\setlength{\parsep}{0pt}
  \item \textbf{TF-IDF}（1--2~gram，max\_features=5,000）：词级稀疏特征，可解释性强，作为经典基线。
  \item \textbf{BioBERT}（\texttt{dmis-lab/biobert-base-cased-v1.1}，mean~pooling至768维）：预训练语言模型，捕捉医学术语的上下文语义关系。
  \item \textbf{LDA}（Latent~Dirichlet~Allocation，$K=15$）：文档级主题分布，捕捉隐式的主题结构。
  \item \textbf{元特征}（出版年份、期刊类型、MeSH存在标志等）：非文本信号，辅助分类决策。
\end{itemize}
多种表示对比可揭示\textbf{语义深度}（TF-IDF$\to$BioBERT）、\textbf{结构粒度}（词级$\to$文档级）和\textbf{信息类型}（文本$\to$非文本）对分类性能的贡献。

\subsection*{2.2 分类算法（12种）}

\vspace{-2pt}
{\footnotesize
\begin{tabular}{@{}lll@{}}
\toprule
\textbf{课程模块} & \textbf{算法} & \textbf{适用子任务} \\
\midrule
贝叶斯学习           & Na{\"\i}ve Bayes & 类别分类 \\
基于实例的学习       & $k$-NN、SVM(RBF) & 类别分类 \\
回归学习             & Logistic Regression & 类别分类 \\
集成学习（Bagging）  & Random Forest & 类别、标签分类 \\
集成学习（Boosting） & AdaBoost, XGBoost & 类别、标签分类 \\
深度学习             & BioBERT+MLP微调 & 类别分类 \\
无监督学习           & LDA主题+聚类 & 文献子领域发现 \\
图表示学习          & node2vec, GCN, GraphSAGE & PGB节点分类 \\
\bottomrule
\end{tabular}\par}
\vspace{2pt}

共\textbf{12种算法}，覆盖课程全部主要模块。图方法仅适用于PGB（天然图结构），其余方法在三个全标注数据集上独立评估。多标签任务额外对比Binary~Relevance（BR，每个标签独立二分类）、Classifier~Chains（CC，链式传递标签依赖）和Label~Powerset（LP，将标签组合转为多分类）三种问题转换策略。

\section{三、数据与实验方案}

\subsection*{3.1 数据集}

\vspace{-2pt}
{\footnotesize
\begin{tabular}{@{}p{2.1cm}p{2.8cm}p{2.5cm}p{2.5cm}p{2.8cm}@{}}
\toprule
\textbf{属性} & \textbf{OHSUMED} & \textbf{PubMed-ML} & \textbf{PGB} & \textbf{Spatial Tracker} \\
\midrule
来源 & MEDLINE 1987--91 & Kaggle & S2ORC+PubMed & PubMed检索(待构建) \\
规模 & $\sim$294K篇 & 10K篇 & $\sim$30M篇 & $\sim$数千篇(待标注) \\
标签空间 & 14,466 MeSH词 & 15粗分类 & 3类+MeSH层级 & 5类+若干标签(待设计) \\
主要任务 & 多标签分类 & 多标签分类 & 节点分类 & 类别+多标签+丢弃 \\
标注状态 & 全量标注 & 全量标注 & 全量标注 & LLM批量标注(待进行) \\
\bottomrule
\end{tabular}\par}
\vspace{2pt}

\subsection*{3.2 渐进式实验设计}

实验按\textbf{三步渐进}组织，每一步以上一步结果为输入：

\textbf{Step~1（多数据集算法筛选）}——在三个全标注数据集上全面对比12种算法和4种文本表示：
\begin{itemize}
  \setlength{\itemsep}{-2pt}\setlength{\parsep}{0pt}
  \item 固定SVM，对比TF-IDF/BioBERT/LDA/组合的5折CV Macro~F1，选出\textbf{最优文本表示}。
  \item 固定BioBERT嵌入，对比12种算法的Macro~F1与训练时间，在OHSUMED（14K稀疏标签）和PubMed-ML（15稠密标签）上分别考察\textbf{标签密度}对算法排序的影响。
  \item 在PGB上，额外引入\textbf{图嵌入+图神经网络}，对比纯文本vs文本+图结构vs文本+图+MeSH层级三档性能，量化结构信息的边际收益。
\end{itemize}

\textbf{Step~2（目标数据集构建与应用）}——构建空间转录组学文献数据集并验证最佳方法：
\begin{itemize}
  \setlength{\itemsep}{-2pt}\setlength{\parsep}{0pt}
  \item 设计PubMed检索式，爬取空间转录组学相关文献，构建带标题+摘要+MeSH词的原始库。
  \item 参考Biomed-Enriched两阶段方法：DeepSeek API标注$\to$蒸馏模型扩标全量$\to$以手动标注子集评估质量。
  \item 将Step~1选出的top-3方法应用于该数据集，对比\textbf{BioBERT+MLP基线}与最优替代方法的性能差异。
\end{itemize}

\textbf{Step~3（迁移微调探索）}——探索\textbf{宽泛$\to$精细}的迁移学习范式：
\begin{itemize}
  \setlength{\itemsep}{-2pt}\setlength{\parsep}{0pt}
  \item 对支持增量学习的算法（BioBERT+MLP、XGBoost warm~start、GCN/GraphSAGE），在OHSUMED或PGB上预训练，在空间转录组学数据上微调，对比微调前后的F1增量。
\end{itemize}

\section{四、评价方案与预期结果}

\subsection*{4.1 评价指标}

类别分类：Accuracy、Macro/Weighted~F1、Cohen's~$\kappa$；多标签：Jaccard、Hamming~Loss、Per-label~F1；图节点分类：Micro/Macro~F1；迁移学习：微调前后F1增量。

\subsection*{4.2 预期结果}

\textbf{(1)~算法排序}：BioBERT+MLP$\ge$XGBoost$\ge$RF$\ge$SVM$\ge$AdaBoost$\ge$LR$\ge$$k$-NN$\ge$NB（OHSUMED上集成方法优势更显著，PubMed-ML上差距缩小）。

\textbf{(2)~图结构价值}：PGB上引入MeSH层级信息和图结构的GCN/GraphSAGE相比纯文本方法预期有一定提升，验证结构信息对文献分类的补充价值。

\textbf{(3)~LLM标注可行性}：DeepSeek标注与人工标注在常见类别上预期有较高一致率，蒸馏后下游F1可接近人工标注水平。

\textbf{(4)~迁移微调增益}：在宽泛生物医学数据上预训练后，于空间转录组学数据上微调预期能进一步提升F1，验证"宽泛预训练+领域微调"范式的有效性。

\end{document}
